\documentclass[11pt,a4paper]{article}

\usepackage[T1]{fontenc}
\usepackage[utf8]{inputenc}
\usepackage{lmodern}
\usepackage{geometry}
\usepackage{setspace}
\usepackage{booktabs}
\usepackage{longtable}
\usepackage{array}
\usepackage{graphicx}
\usepackage{xcolor}
\usepackage{hyperref}
\usepackage{enumitem}
\usepackage{float}
\usepackage{amsmath}
\usepackage{caption}

\geometry{margin=2.5cm}
\setstretch{1.15}
\hypersetup{
    colorlinks=true,
    linkcolor=blue,
    urlcolor=blue,
    pdftitle={Relatorio do Dataset de Grupo - Analise de Dados e Modelacao em KNIME}
}

\newcommand{\placeholderfig}[2]{
\begin{figure}[H]
    \centering
    \fbox{%
        \begin{minipage}[c][6cm][c]{0.9\textwidth}
            \centering
            \textbf{[INSERIR IMAGEM KNIME]}
        \end{minipage}}
    \caption{#2}
\end{figure}
}

\title{\textbf{Relatorio de Resultados}\\
Dataset de Grupo: Spotify Tracks\\
\large Analise de Dados e Modelacao em KNIME}
\author{\textbf{Unidade Curricular:} Analise e Descoberta de Informacao \\
\textbf{Grupo:} [Preencher] \\
\textbf{Elementos:} [Preencher]}
\date{\today}

\begin{document}

\maketitle
\thispagestyle{empty}
\clearpage
\tableofcontents
\clearpage
\pagenumbering{arabic}

\section{Introducao}
Este relatorio documenta, de forma unica e coerente, o trabalho desenvolvido sobre o \textbf{dataset de grupo} selecionado no projeto, com implementacao em \textit{workflows} KNIME e enquadramento no metodo CRISP-DM. O objetivo nao foi apenas construir modelos, mas sim mostrar todo o percurso analitico: compreensao do problema, compreensao dos dados, preparacao, modelacao, avaliacao e interpretacao critica dos resultados.

O problema analisado centra-se num conjunto de faixas musicais do ecossistema Spotify, com metadados editoriais, atributos acusticos e uma variavel alvo de interesse ligada a popularidade. A partir deste contexto, foram exploradas duas linhas de trabalho complementares. A primeira procura agrupar faixas em perfis sonoros semelhantes, usando \textit{clustering} e analise de dimensionalidade. A segunda pretende prever a popularidade das faixas por regressao, comparando modelos lineares e nao lineares. Em paralelo, o workflow inclui ainda uma experiencia auxiliar de classificacao por rede neuronal, usada como verificacao adicional da separabilidade dos grupos obtidos.

O texto foi escrito para suportar a reproducao do trabalho e para refletir o que esta realmente implementado no workflow, evitando descricoes genericas ou passos que nao estejam presentes no projeto.

\section{Dominio, objetivos e abordagem}
\subsection{Dominio analisado}
O dominio escolhido e o da analise musical e de consumo digital. Neste contexto, as caracteristicas acusticas das faixas podem ser usadas para segmentacao editorial, recomendacao e estudo de perfis sonoros. Variaveis como \texttt{danceability}, \texttt{energy}, \texttt{loudness}, \texttt{acousticness} e \texttt{tempo} capturam propriedades com interpretacao musical direta, enquanto \texttt{genre}, \texttt{release\_year} e \texttt{explicit} acrescentam contexto de catalogo.

\subsection{Objetivos}
Os objetivos do trabalho podem resumir-se em quatro pontos. Primeiro, analisar a qualidade do dataset e demonstrar os problemas existentes antes da modelacao. Segundo, construir um pipeline de preparacao robusto, capaz de remover duplicados, tratar valores em falta, controlar outliers e criar variaveis derivadas. Terceiro, conceber e comparar varios modelos de machine learning adequados ao problema de segmentacao e previsao. Quarto, produzir uma leitura critica dos resultados, identificando o que o workflow ja demonstra com solidez e o que ainda poderia ser aprofundado.

\subsection{Estrategia seguida}
A estrategia adotada foi deliberadamente iterativa. A leitura inicial do dataset orientou a preparacao dos dados; a preparacao definiu quais as features disponiveis para os modelos; e a modelacao devolveu evidencia que voltou a influenciar a interpretacao do dataset. Esta forma de trabalho e totalmente consistente com CRISP-DM, porque permite voltar atras sempre que a analise revela problemas de qualidade ou de interpretacao.

\section{Metodologia de analise de dados}
Neste projeto foi aplicado o metodo \textbf{CRISP-DM} (Cross-Industry Standard Process for Data Mining). O seu valor aqui e sobretudo pratico: cada fase responde a uma pergunta concreta e gera artefactos que alimentam a fase seguinte.

Na fase de \textbf{compreensao do problema}, foi definida a finalidade analitica do dataset de grupo: identificar perfis musicais coerentes e estimar a popularidade das faixas. Na fase de \textbf{compreensao dos dados}, o workflow confirmou a estrutura do CSV, mediu a distribuicao das variaveis e expôs problemas como faltas, outliers, intervalos de dominio violados e duplicados completos.

A fase de \textbf{preparacao} foi a mais rica tecnicamente. Em vez de aplicar uma limpeza superficial, o workflow faz uma sanitizacao controlada: primeiro remove duplicados, depois converte e corrige campos invalidos, em seguida imputa valores em falta e so depois trata extremos numericos. Este encadeamento e importante porque evita que os modelos aprendam com ruido desnecessario e permite manter a maior quantidade possivel de informacao util.

Na fase de \textbf{modelacao}, foram construidos ramos distintos para clustering e regressao. O clustering foi avaliado com o metodo do cotovelo, silhouette e PCA, enquanto a regressao foi suportada por treino/teste estratificado, normalizacao e comparacao entre regressao linear, arvore de regressao e um ramo de random forest ainda por executar.

Na fase de \textbf{avaliacao}, o foco nao foi apenas a qualidade numerica dos modelos, mas tambem a interpretacao do seu significado: se os clusters sao realmente separaveis, se a regressao linear e suficiente ou se existem evidencias de nao linearidade, e quao robusta e a preparacao perante dados fora de dominio.

\section{Descricao do dataset}
\subsection{Estrutura geral}
O ficheiro \texttt{spotify\_tracks.csv} contem \textbf{100\,500 registos} e \textbf{21 colunas}. O periodo temporal abrangido vai de \textbf{2000 a 2024} e o dataset cobre \textbf{20 generos} musicais. Existe ainda um ficheiro complementar, \texttt{spotify\_tracks\_errors\_log.csv}, que regista os erros introduzidos intencionalmente no processo de geracao do dataset, o que foi essencial para validar a preparacao de dados.

\begin{table}[H]
\centering
\caption{Estrutura do dataset de grupo}
\begin{tabular}{p{4cm}p{10cm}}
\toprule
Grupo de variaveis & Exemplos \\
\midrule
Identificacao & \texttt{track\_id}, \texttt{track\_name}, \texttt{artist\_name}, \texttt{album\_name} \\
Contexto & \texttt{release\_year}, \texttt{genre}, \texttt{explicit}, \texttt{key}, \texttt{mode}, \texttt{time\_signature} \\
Target & \texttt{popularity} \\
Audio & \texttt{danceability}, \texttt{energy}, \texttt{loudness}, \texttt{speechiness}, \texttt{acousticness}, \texttt{instrumentalness}, \texttt{liveness}, \texttt{valence}, \texttt{tempo}, \texttt{duration\_ms} \\
\bottomrule
\end{tabular}
\end{table}

\subsection{Qualidade dos dados}
O dataset nao foi criado para ser limpo; foi criado para simular um ambiente realista de preparacao. O gerador introduz valores em falta, outliers, violacoes de dominio e duplicados. Isto significa que a analise nao podia começar diretamente pela modelacao. Era necessario demonstrar, com evidencia, que os problemas estavam identificados e que o workflow sabia resolvê-los de forma consistente.

Em termos praticos, os problemas mais relevantes sao os seguintes. Existem valores em falta nas features de audio, o que afeta tanto a analise exploratoria como a modelacao. Existem linhas duplicadas completas, que distorcem contagens, distribuicoes e treino de modelos. Existem ainda valores fora do dominio esperado em variaveis limitadas a intervalos bem definidos, como as features normalizadas entre 0 e 1, e outliers fortes em variaveis como \texttt{tempo}, \texttt{loudness} e \texttt{duration\_ms}.

\subsection{Exploracao inicial}
A fase de exploracao foi organizada em torno de varios blocos visuais no KNIME. O objetivo era verificar rapidamente se o dataset estava lido corretamente, perceber a forma das distribuicoes e encontrar sinais de problemas a tratar.

\placeholderfig{Distribuicao da popularidade e das principais features numericas no KNIME.}{Exploracao univariada do dataset Spotify}
\placeholderfig{Boxplots para deteccao visual de outliers e violacoes de dominio.}{Outliers e dispersao inicial}
\placeholderfig{Scatter plots entre variaveis sonoras para identificar relacoes entre atributos.}{Relacoes bivariadas entre features acusticas}
\placeholderfig{Distribuicao por genero musical para perceber o equilibrio do dataset.}{Representatividade por genero}

A exploracao mostrou que as variaveis acusticas tinham escalas muito diferentes, o que e natural, mas tambem confirmou que havia varios pontos de anomalia a tratar. Por exemplo, alguns valores de \texttt{danceability}, \texttt{valence} ou \texttt{energy} surgem fora do intervalo esperado. Em paralelo, o comportamento de \texttt{tempo} e \texttt{loudness} inclui valores extremos que justificam tratamento especifico. Esta leitura inicial explica porque o workflow nao passa diretamente da leitura do CSV para a modelacao.

\section{Preparacao dos dados em KNIME}
A preparacao foi desenhada como um pipeline sequencial, onde cada no acrescenta uma operacao concreta ao estado dos dados.

\subsection{Remocao de duplicados}
O primeiro passo foi o \texttt{Duplicate Row Filter}. A escolha de tratar duplicados antes de qualquer outra transformacao evita que a limpeza ou a imputacao sejam aplicadas duas vezes aos mesmos registos. O workflow compara praticamente todas as colunas relevantes e conserva a primeira ocorrencia, mantendo a ordem original. O resultado observado e claro: de 100\,500 linhas iniciais, o dataset passa para 100\,000 linhas, o que confirma a existencia de 500 duplicados removidos no inicio do pipeline.

\subsection{Sanitizacao e engenharia de atributos}
Depois dos duplicados, o no \texttt{Expression} concentra varias operacoes que poderiam estar espalhadas por mais nodos. Aqui sao feitas tres coisas muito importantes. Primeiro, valores invalidos em variaveis limitadas a dominos conhecidos sao convertidos em missing para serem resolvidos mais tarde com regra consistente. Segundo, a variavel \texttt{explicit} e transformada para uma representacao numerica binaria. Terceiro, sao criadas features derivadas como \texttt{duration\_min} e a classe auxiliar \texttt{popularity\_class}.

Esta etapa e crucial porque transforma o dataset num formato mais adequado a aprendizagem automatica sem perder a informacao original de contexto. A criacao de \texttt{duration\_min} torna a duracao mais interpretabvel do que \texttt{duration\_ms}, e a variavel \texttt{popularity\_class} permite estratificar o split de treino/teste e apoiar a leitura dos resultados.

\subsection{Tratamento de valores em falta}
O no \texttt{Missing Value} recebe tanto os missing originais do CSV como os que sao gerados pelo passo anterior. A estrategia nao e uniforme; e precisamente essa diferenca que mostra maturidade metodologica. Para algumas features, o workflow usa media; para outras, mediana; e para o target \texttt{popularity}, a linha e removida quando o valor esta em falta.

\begin{table}[H]
\centering
\caption{Estrategias de imputacao no no Missing Value}
\begin{tabular}{p{4cm}p{8cm}}
\toprule
Coluna & Estrategia \\
\midrule
\texttt{danceability} & Mean \\
\texttt{valence} & Mean \\
\texttt{energy}, \texttt{speechiness}, \texttt{acousticness}, \texttt{instrumentalness}, \\
\texttt{liveness}, \texttt{loudness}, \texttt{tempo}, \texttt{duration\_ms} & Median \\
\texttt{popularity} & Remove Row \\
\bottomrule
\end{tabular}
\end{table}

Depois desta fase, o dataset fica com 99\,353 linhas. O valor e importante porque confirma que a preparacao nao removeu tudo indiscriminadamente; removeu apenas o que era metodologicamente necessario remover.

\subsection{Tratamento de outliers}
Os outliers continuos sao tratados num no especifico, \texttt{Numeric Outliers}, que opera sobre \texttt{duration\_ms}, \texttt{loudness} e \texttt{tempo}. O criterio de deteccao e o IQR, com scalar 3.0, e a acao escolhida e substituir o valor fora de intervalo pelo limite admissivel mais proximo. Esta decisao e relevante porque preserva a linha e impede que observacoes extremas deformem a modelacao.

\subsection{Selecao de variaveis e separacao de ramos}
Depois da limpeza principal, o workflow organiza uma base comum com 15 colunas. Nela permanecem as variaveis que sao uteis para clustering e regressao, incluindo \texttt{release\_year}, \texttt{explicit}, \texttt{danceability}, \texttt{energy} e \texttt{loudness}. Tambem sao mantidas \texttt{tempo}, \texttt{duration\_min} e a classe auxiliar \texttt{popularity\_class}. Ficam de fora identificadores textuais e metadados menos uteis para modelacao direta.

A partir daqui o workflow divide-se em dois ramos. No ramo de clustering, o espaco e reduzido para 10 features sonoras e depois normalizado por Min-Max para o intervalo [0,1]. No ramo de regressao, o dataset e separado em treino e teste com split estratificado de 80/20 sobre \texttt{popularity\_class}, e ambos os subconjuntos sao normalizados de forma consistente. O uso de \texttt{Table Partitioner} com seed fixa garante reprodutibilidade.

\placeholderfig{Fluxo de preparacao com remocao de duplicados, imputacao e outliers.}{Pipeline de preparacao no KNIME}
\placeholderfig{Normalizacao e separacao de treino/teste para a regressao.}{Preparacao final para modelacao}

\section{Modelacao}
\subsection{Clustering}
O clustering procura agrupar faixas com perfis sonoros semelhantes. Para isso o workflow usa apenas as features acusticas normalizadas, o que e a opcao correta porque o objetivo nao e agrupar por ano, genero ou identificador, mas sim por assinatura sonora.

\subsubsection*{Metodo do cotovelo}
A escolha do numero de clusters foi analisada com varias execucoes de k-Means para valores de \textit{k} entre 3 e 9. O workflow nao usa uma funcao automatica pronta para o elbow; em vez disso, reconstrói manualmente a soma das distancias quadradas dentro de cada cluster e agrega os resultados num grafico de linha. Esse detalhe e importante porque mostra que a decisao de \textit{k} nao foi arbitraria: foi suportada por uma curva de WCSS produzida no proprio KNIME.

\placeholderfig{Curva do elbow para suportar a escolha do numero de clusters.}{Metodo do cotovelo}

\subsubsection*{k-Means final}
O modelo final de clustering foi fixado em \textbf{5 clusters}, com inicializacao aleatoria, seed 1 e maximo de 99 iteracoes. O output do proprio no inclui a coluna de cluster, pelo que nao foi necessario um assigner separado. Esta decisao final e coerente com a leitura da curva do elbow e com a interpretacao posterior dos centroides.

\placeholderfig{Visualizacao em PCA dos pontos agrupados por cluster.}{Clusters projetados em duas dimensoes}
\placeholderfig{Perfil dos centroides por feature para interpretar cada cluster.}{Centroides dos clusters}

\subsubsection*{Interpretacao dos clusters}
Para tornar os clusters interpretaveis, o workflow devolve os centroides a forma longa, cria graficos de barras e junta o genero musical para verificar composicao de cada grupo. Esta parte e essencial porque um cluster sem leitura semantica tem pouco valor analitico. Aqui a interpretacao passa a ser musical: clusters mais energeticos e dançaveis podem representar uma classe de faixas com uso mais intenso de ritmo, enquanto clusters com maior acustica e menor energia podem corresponder a perfis mais calmos ou instrumentais.

A qualidade do agrupamento foi analisada adicionalmente com silhouette. O metodo hierarquico sobre uma amostra obteve um valor medio global de aproximadamente \textbf{0.0732}, o que e baixo e sugere sobreposicao entre grupos. Em ramos de amostragem estratificada por cluster, surgem valores de cerca de \textbf{0.1609}, \textbf{0.1782} e \textbf{0.1740}, que continuam modestos mas sao superiores ao primeiro valor observado. Isto mostra que o clustering tem utilidade exploratoria, mas nao apresenta separacao perfeita entre classes sonoras.

\placeholderfig{Resultado de silhouette para a comparacao entre solucoes de clustering.}{Silhouette e qualidade do agrupamento}
\placeholderfig{Composicao dos clusters por genero musical.}{Distribuicao dos generos por cluster}

\subsubsection*{Comparacao hierarquica}
O workflow inclui ainda um ramo de clustering hierarquico com distancia euclidiana e linkage complete. Este ramo nao substitui o k-Means final, mas serve para comparar a qualidade da segmentacao. A leitura critica deste valor de silhouette e importante: o objetivo nao e apenas escolher um algoritmo, mas perceber se a estrutura dos dados suporta clusters nitidos ou se existe uma sobreposicao natural entre estilos musicais.

\subsection{Regressao}
O segundo objetivo de modelacao e prever \texttt{popularity}. Ao contrario do clustering, aqui ha uma variavel target numerica e, por isso, a interpretacao dos resultados passa pela qualidade da previsao e nao pela separacao entre grupos.

\subsubsection*{Regressao linear}
O primeiro baseline e uma regressao linear multipla. O learner usa \texttt{popularity} como target e inclui constante. A vantagem deste modelo e a interpretabilidade: se o comportamento da popularidade for aproximadamente linear, este baseline pode ser suficiente. O workflow inclui ainda grafico de dispersao entre valores reais e previstos e histogramas de residuos para verificar se os erros estao centrados e se ha padroes sistematicos nao capturados.

\placeholderfig{Comparacao entre popularidade real e prevista pela regressao linear.}{Regressao linear: real vs previsto}
\placeholderfig{Distribuicao dos residuos da regressao linear.}{Residuos da regressao linear}

\subsubsection*{Arvore de regressao}
Como alternativa nao linear, o workflow inclui uma simple regression tree. Os limites da arvore estao praticamente abertos, sem poda agressiva, o que lhe permite capturar limiares e interacoes locais. Este modelo e particularmente util quando a popularidade depende de zonas especificas do espaco de features, e nao apenas de efeitos globais lineares.

\placeholderfig{Comparacao entre regressao linear e arvore de regressao.}{Comparacao de modelos de regressao}
\placeholderfig{Residuos da arvore de regressao.}{Analise de residuos da arvore}

\subsubsection*{Random forest de regressao}
Existe ainda um ramo de random forest de regressao configurado com 100 arvores, seed 1 e bootstrap ativo, mas nao executado na versao atual do workflow. Isto nao invalida o relatorio; pelo contrario, mostra que a comparacao de modelos esta montada para evoluir. Do ponto de vista analitico, o ramo sugere uma proximidade metodologica a uma comparacao mais rica entre modelos lineares e ensembles, que seria o passo seguinte natural.

\subsubsection*{Experiencia auxiliar de classificacao}
O workflow inclui ainda uma branch auxiliar de classificacao com RProp MLP para prever o cluster. Esta experiencia nao e o foco principal do projeto, mas tem um papel util: verifica se os clusters obtidos sao suficientemente separados para serem aprendidos por um classificador. O modelo atingiu \textbf{accuracy de 0.9652418976}, com \textbf{Cohen's kappa de 0.9538043842}, o que indica uma separabilidade elevada dos grupos obtidos pelo clustering.

\section{Resultados e analise critica}
\subsection{Resumo dos resultados mais relevantes}
\begin{table}[H]
\centering
\caption{Resumo dos principais resultados observados}
\begin{tabular}{p{5cm}p{8.5cm}}
\toprule
Elemento & Resultado observado \\
\midrule
Tamanho inicial do dataset & 100\,500 registos e 21 colunas \\
Duplicados removidos & 500 linhas \\
Base comum apos preparacao & 99\,353 linhas e 15 colunas \\
Ramo de clustering & 99\,353 linhas e 10 features normalizadas \\
Split treino/teste & 80/20, estratificado por \texttt{popularity\_class}, seed 1 \\
Silhouette hierarquico & aproximadamente 0.0732 \\
Silhouettes adicionais & aproximadamente 0.1609, 0.1782 e 0.1740 \\
Cluster final & \textit{k} = 5 \\
Classificador auxiliar do cluster & accuracy de 0.9652418976 \\
\bottomrule
\end{tabular}
\end{table}

\subsection{Leitura critica do clustering}
Os resultados de clustering mostram que existe estrutura nos dados, mas a separacao nao e perfeita. Um silhouette perto de 0.07 e um sinal de clusters relativamente sobrepostos, o que faz sentido num dataset musical: as faixas podem partilhar intensidades semelhantes de energia, valencia e dançabilidade, sem que isso se traduza em grupos radicalmente distintos. Em termos práticos, isto significa que o clustering e mais util como ferramenta de segmentacao e exploracao do que como classificacao rigorosa e definitiva.

Ao mesmo tempo, a forte accuracy do classificador auxiliar sugere que, apesar da sobreposicao observada no silhouette, os clusters sao suficientemente estruturados para serem reaprendidos por um modelo supervisionado. Esta combinacao de sinais e interessante: os clusters nao sao perfeitamente separados no espaco geometrico, mas ainda assim mantem uma identidade suficiente para uso operacional.

\subsection{Leitura critica da regressao}
Na regressao, a interpretacao deve ser feita com mais cautela porque o workflow nao fecha ainda uma tabela consolidada de MAE, RMSE e \(R^2\) para todos os modelos. O que o workflow mostra com seguranca e que a preparacao foi feita sem fuga de informacao, a regressao linear e a arvore estao executadas e os residuos foram diagnosticados visualmente.

Do ponto de vista metodologico, a regressao linear funciona como baseline interpretavel. A arvore de regressao tem potencial para capturar nao linearidades, mas a ausencia de poda rigorosa pode aumentar a variancia. O ramo de random forest, embora ainda nao executado, e precisamente o tipo de modelo que pode oferecer um compromisso mais robusto entre viés e variancia. Por isso, a comparacao futura entre estes tres modelos e o passo natural seguinte.

\subsection{Limitações do estudo}
A principal limitacao e que o dataset e sintético, ainda que inspirado em distribuicoes realistas. Isso e util para treino, mas nao substitui completamente um catalogo real com variacoes de metadata mais complexas. Outra limitacao e a ausencia, no estado atual do workflow, de uma tabela unica consolidada com metrias de regressao finalizadas. Por fim, o trabalho analisa associacoes e padroes, mas nao pode afirmar causalidade entre as features e a popularidade.

\placeholderfig{Erro de regressao e padrões residuais observados no workflow.}{Analise dos residuos e limite dos modelos}

\section{Sugestoes e recomendacoes}
A primeira recomendacao e completar a comparacao de regressao com a execucao do ramo de random forest e a consolidacao das metricas finais numa tabela unica. A segunda e testar diferentes numeros de clusters ou diferentes seeds para verificar estabilidade da segmentacao. A terceira e aprofundar a interpretacao dos clusters por genero, por exemplo associando centroides a perfis musicais mais claros. A quarta e explorar importancia das variaveis em regressao para identificar quais as features acusticas que mais contribuem para popularidade.

Do ponto de vista pratico, o workflow poderia ainda ganhar robustez com validacao cruzada em regressao, repetição de experiencias com seeds diferentes e comparação entre escalas de normalizacao. Essas extensões não são necessárias para concluir o relatório, mas seriam um bom passo seguinte se o projeto continuasse.

\section{Reprodutibilidade}
Para reproduzir o trabalho, basta seguir a sequencia geral do workflow. Primeiro, abrir o ficheiro do KNIME relativo ao dataset de grupo. Depois, executar o bloco de \textit{data understanding} para confirmar leitura, distribuicoes e duplicados. Em seguida, correr o pipeline de preparacao na ordem definida: remocao de duplicados, expression, missing value, numeric outliers e filtragem final. Depois, correr separadamente o ramo de clustering e o ramo de regressao, garantindo que a normalizacao e a particao sao executadas antes dos modelos.

Os ficheiros que sustentam a reproducao sao o CSV principal, o log de erros, o script de geracao e o workflow KNIME. No caso de querer repetir a analise de raiz, convem confirmar se os caminhos dos ficheiros estao corretos e se as ligacoes no KNIME apontam para a pasta do projeto.

\section{Conclusao}
O trabalho desenvolvido sobre o dataset de grupo mostra uma cadeia completa de analise de dados: exploracao, preparacao, modelacao e interpretacao. A maior forca do projeto esta na qualidade da preparacao e na coerencia entre o que foi visto nos dados e o que foi feito no workflow. No clustering, o resultado final e suficientemente rico para segmentacao e interpretacao musical. Na regressao, o pipeline esta bem montado e tem comparacao real entre abordagens, embora ainda beneficie de consolidacao de metricas.

Em suma, o relatorio documenta um processo analitico serio, reproduzivel e alinhado com CRISP-DM, com suficiente detalhe para justificar as escolhas feitas e com espaco claro para evolucao futura do modelo.

\end{document}
